\documentclass[
	12pt,				% tamanho da fonte
	openany,			% capítulos começam em pág ímpar (insere página vazia caso preciso)
	oneside, 			% oneside - twoside
	a4paper,			% tamanho do papel.
	chapter=TITLE,		% títulos de capítulos convertidos em letras maiúsculas
	section=TITLE,		% títulos de seções convertidos em letras maiúsculas
	sumario=tradicional,	
	%subsection=TITLE,	% títulos de subseções convertidos em letras maiúsculas
	%subsubsection=TITLE,% títulos de subsubseções convertidos em letras maiúsculas
	english,			% idioma adicional para hifenização
	brazil,				% o último idioma é o principal do documento
	]{abntex2}

% ---------------------------------------------------------------------------
% Inclui os comandos do projeto
% ---------------------------------------------------------------------------
% -----------------------------------------------------------------------------
% Pacotes fundamentais
% -----------------------------------------------------------------------------
\usepackage{xcolor}
\newcommand\myworries[1]{\textcolor{red}{[#1]}}
% Escolhendo a fonte
% Consultar: http://www.tug.dk/FontCatalogue/
%\usepackage{lmodern}       % Usa a fonte Latin Modern (Serifada, tipo Times New Roman
\usepackage{fourier}        % Usa a fonte Utopia Regular with Fourier
\usepackage[T1]{fontenc}    % ^^^^

%\usepackage{helvet}		% Usa a fonte Helvetica (Tipo Arial)	
%\renewcommand{\familydefault}{\sfdefault} tira o serifado
\usepackage[T1]{fontenc}		% Selecao de codigos de fonte.
\usepackage[utf8]{inputenc}		% Codificacao do documento (conversão automática dos acentos)
\usepackage{indentfirst}		% Indenta o primeiro parágrafo de cada seção.
\usepackage{color}				% Controle das cores
\usepackage{tikz}				% Inclusão de gráficos
\usepackage{graphicx}			% Inclusão de gráficos
\usepackage{microtype} 			% para melhorias de justificação
% -----------------------------------------------------------------------------
% Pacotes adicionais, usados no anexo do modelo de folha de identificação
% -----------------------------------------------------------------------------
\usepackage{multicol}
\usepackage{multirow}
% -----------------------------------------------------------------------------
% Pacotes adicionais, usados apenas no âmbito do Modelo Canônico do abnteX2
% -----------------------------------------------------------------------------
\usepackage{lipsum}				% para geração de dummy text % http://ctan.org/pkg/lipsum
% -----------------------------------------------------------------------------
% Pacotes de citações
% -----------------------------------------------------------------------------
\usepackage[brazilian,hyperpageref]{backref}	 % Paginas com as citações na bibliografia
\usepackage[alf,abnt-etal-list=3,abnt-etal-cite=3, abnt-emphasize=bf,num]{abntex2cite}	% Citações padrão ABNT
\usepackage{pdflscape}
\usepackage{footnote}
\usepackage{pdfpages}
\usepackage{caption}

% -----------------------------------------------------------------------------
% Pacotes adicionados por @leolleocomp
% ----------------------------------------------------------------------------- 
\usepackage{booktabs}
\usepackage{adjustbox}
\usepackage{subcaption}
\usepackage[labelfont=bf]{caption}
\usepackage{gensymb}
\usepackage{amsmath}
\usepackage{array}
\usepackage{float}
\usepackage{xcolor,colortbl}
\usepackage{longtable}
\usepackage{scalefnt}
\usepackage{listings}			% inserir codigo fonte


% -----------------------------------------------------------------------------
% Pacotes adicionados por @Gabrielr2508
% ----------------------------------------------------------------------------- 
\usepackage{hyperref}
\PassOptionsToPackage{hyphens}{url}\usepackage{hyperref}

\usepackage{tocloft}
% -- permite a adição de células especiais em tabelas
\newcommand{\specialcell}[2][c]{%
  \begin{tabular}[#1]{@{}c@{}}#2\end{tabular}}

\newcounter{equationset}
\newcommand{\equationset}[1]{% \equationset{<caption>}
  \refstepcounter{equationset}% Step counter
  \noindent\makebox[\linewidth]{Equação ~\theequationset: #1}
 }

%--------------------------------------------------------------------------------
% Adequação dos títulos dos capitulos, seções, subseções às normas da Univasf
% Added by @Gabrielr2508
%--------------------------------------------------------------------------------
\renewcommand{\ABNTEXchapterfont}{\fontseries{b}}
\renewcommand{\ABNTEXchapterfontsize}{\normalsize}

\renewcommand{\ABNTEXsectionfont}{\fontseries{m}}
\renewcommand{\ABNTEXsectionfontsize}{\normalsize}

\renewcommand{\ABNTEXsubsectionfont}{\fontseries{b}}
\renewcommand{\ABNTEXsubsectionfontsize}{\normalsize}

\renewcommand{\ABNTEXsubsubsectionfont}{\fontseries{m}}
\renewcommand{\ABNTEXsubsubsectionfontsize}{\normalsize}

%--------------------------------------------------------------------------------
% CONFIGURAÇÕES DE PACOTES
% Configurações do pacote backref
%--------------------------------------------------------------------------------
% Usado sem a opção hyperpageref de backref
% \renewcommand{\backrefpagesname}{Citado na(s) página(s):~}
% Texto padrão antes do número das páginas
\renewcommand{\backref}{}
% Define os textos da citação
% \renewcommand*{\backrefalt}[4]{
% 	\ifcase #1 %
% 		%Nenhuma citação no texto.%
% 	\or
% 		Citado na página #2.%
% 	\else
% 		Citado #1 vezes nas páginas #2.%
% 	\fi}%

%--------------------------------------------------------------------------------
% Configurações de aparência do PDF final
%--------------------------------------------------------------------------------
% alterando o aspecto da cor azul
\definecolor{blue}{RGB}{41,5,195}

% informações do PDF
\makeatletter
\hypersetup{
     	%pagebackref=true,
		pdftitle={\@title},
		pdfauthor={\@author},
    	pdfsubject={\imprimirpreambulo},
	    pdfcreator={LaTeX with abnTeX2},
		pdfkeywords={abnt}{latex}{abntex}{abntex2}{relatório técnico},
		colorlinks=true,			% false: boxed links; true: colored links
    	linkcolor=black,				% color of internal links
    	citecolor=black,				% color of links to bibliography
    	filecolor=black,			% color of file links
		urlcolor=black,
		bookmarksdepth=4
}
\makeatother
% ---

% ---
% Espaçamentos entre linhas e parágrafos
% ---

% O tamanho do parágrafo é dado por:
\setlength{\parindent}{1.3cm}

% Controle do espaçamento entre um parágrafo e outro:
\setlength{\parskip}{0.2cm}  % tente também \onelineskip

%--------------------------------------------------------------------------------
% compila o indice
%--------------------------------------------------------------------------------
\makeindex
% ---

%--------------------------------------------------------------------------------
% Comando para inserir imagens de forma simples
%--------------------------------------------------------------------------------
\newcommand{\imagem}[4]
{%			\imagem{x.x}{nomeimg}{titulo}{fonte}
	\begin{figure}[!htb]
		\caption{\label{img:#2}#3}
		\begin{center}
			\includegraphics[scale=#1]{img/#2}
		\end{center}
        \legend{\textbf{Fonte:} #4}
	\end{figure}
}%

%--------------------------------------------------------------------------------
% Creio que esses comandos sejam para desenhar algo, aguardando explicações de @leolleocomp
%--------------------------------------------------------------------------------
\newcommand{\xx} {$\bigotimes$}
\newcommand{\oo} {$\bigcirc$}

%--------------------------------------------------------------------------------
% Biblioteca para códigos-fonte
%--------------------------------------------------------------------------------
\usepackage[newfloat=true]{minted}

%--------------------------------------------------------------------------------
% Caixas batutas - by @leolleocomp
%--------------------------------------------------------------------------------
\usepackage[most]{tcolorbox}
\tcbuselibrary{breakable}

\tcbuselibrary{minted}
\tcbset{listing engine=minted}

\definecolor{bg}{rgb}{0.95,0.95,0.95}

\SetupFloatingEnvironment{listing}{name=Código, listname=Lista de códigos}

%--------------------------------------------------------------------------------
% configuração do contador dos códigos-fonte - by @leolleocomp
% assim como as figuras, começa em 1
\newcounter{sourcecode}
%--------------------------------------------------------------------------------
%--------------------------------------------------------------------------------
% @leolleocomp
% stackoverflow code
% peguei da resposta abaixo
% https://stackoverflow.com/questions/24086366/change-latex-minted-listings-numbering-to-include-current-section?answertab=votes#tab-top
%--------------------------------------------------------------------------------
\makeatletter
\renewcommand*{\thelisting}{\thesourcecode}
\makeatother

%--------------------------------------------------------------------------------
% Peçam explicações a @leolleo
% WHO DID THIS?
%--------------------------------------------------------------------------------
\newcommand{\Ididthis}{
%	\legend{\textbf{Fonte:} O autor (\the\year).}
\legend{\textbf{Fonte:} O autor}
}

\newcommand{\Otherguydidthis}[1]{
	\legend{\textbf{Fonte:} \citeonline{#1}.}
}

%--------------------------------------------------------------------------------
% Comando para inserir códigos - by @leolleocomp
%--------------------------------------------------------------------------------
\newcommand{\sourcecode}[4]{
\begin{listing}[H]
	\refstepcounter{sourcecode}
	\caption{#1}
	\label{cmd:#2}
	\inputminted[linenos, bgcolor=bg, tabsize=4,breaklines]{#3}{codes/#4}
	\Ididthis
\end{listing}
}


% -----------------------------------------------------------------------------
% Pacotes adicionados por @ruanmed
% ----------------------------------------------------------------------------- 
\usepackage[binary-units=true]{siunitx}


%--------------------------------------------------------------------------------
% Comando para inserir códigos - by @ruanmed
%--------------------------------------------------------------------------------
\usepackage{caption}

\newenvironment{code}{\captionsetup{type=listing}}{}
\SetupFloatingEnvironment{listing}{name=Código}

\newcommand{\sourcecodenolist}[4]{
	\begin{code}
	    \refstepcounter{sourcecode}
        \captionof{listing}{#1 }
        \label{code:#2}
        \inputminted[linenos, bgcolor=bg, tabsize=4,breaklines]{#3}{codes/#4}
        \Wedidthis
    \end{code}
}

% \newcommand{\sourcecodenolist}[4]{
% 	\refstepcounter{sourcecode}
% 	\captionof{listing}{#1 \label{cmd:#2}}
% 	\inputminted[linenos, bgcolor=bg, tabsize=4,breaklines]{#3}{codes/#4}
% 	\Wedidthis
% }

% \newcommand{\sourcecodenolist}[4]{
% 	\refstepcounter{sourcecode}
% 	\captionof{listing}{#1 \label{cmd:#2}}
% 	\inputminted[linenos, bgcolor=bg, tabsize=4,breaklines]{#3}{codes/#4}
% 	\Wedidthis
% }


\newcommand{\sourcecodeinline}[2]{
	\mintinline[linenos, bgcolor=bg, tabsize=4,breaklines]{#1}{#2}
}

%--------------------------------------------------------------------------------
% Outros pacotes adicionais para circuitos - by @ruanmed
%--------------------------------------------------------------------------------
\RequirePackage{tikz}
\usepackage{tikz}
\usepackage[siunitx]{circuitikz}			% para habilitar o desenho de circuitos
\usetikzlibrary{babel}
% ---


%--------------------------------------------------------------------------------
% Outros pacotes adicionais para Siglas - by @ruanmed
%--------------------------------------------------------------------------------
\usepackage{nomencl}
\makenomenclature

\newlist{inlinelist}{enumerate*}{1}
\setlist*[inlinelist,1]{%
  label=(\alph*),
}

%--------------------------------------------------------------------------------
% Pacote para tabelas
%--------------------------------------------------------------------------------
% Pacote para inserir divisória diagonal em tabela
\usepackage{diagbox}

% Pacote para auxiliar em tabelas criando novas células
\usepackage{makecell}

\usepackage{pbox,ragged2e}

\usepackage{graphicx}

%--------------------------------------------------------------------------------
% Outros pacotes adicionais para Caixas de TExto - by @ruanmed
% https://tex.stackexchange.com/questions/13506/how-to-continue-the-framed-text-box-on-multiple-pages
% https://tex.stackexchange.com/questions/34088/how-to-draw-frame-with-rounded-corners-around-box
% https://tex.stackexchange.com/questions/26002/fit-text-into-given-box-by-adjusting-the-fontsize
%--------------------------------------------------------------------------------

\usepackage[most]{tcolorbox}

\usepackage[framemethod=TikZ]{mdframed}

\usepackage{environ}% http://ctan.org/pkg/environ
\newdimen\fontdim
\newdimen\upperfontdim
\newdimen\lowerfontdim
\newif\ifmoreiterations
\fontdim12pt

\makeatletter
\NewEnviron{fitbox}[2]{% \begin{fitbox}{<width>}{<height>} stuff \end{fitbox}
  \def\buildbox{%
    \setbox0\vbox{\hbox{\minipage{#1}%
      \fontsize{\fontdim}{1.2\fontdim}%
      \selectfont%
      \stuff%
    \endminipage}}%
    \dimen@\ht0
    \advance\dimen@\dp0
  }
  \def\stuff{\BODY}% Store environment body
  \buildbox
  % Compute upper and lower bounds
  \ifdim\dimen@>#2
    \loop
      \fontdim.5\fontdim % Reduce font size by half
      \buildbox
    \ifdim\dimen@>#2 \repeat
    \lowerfontdim\fontdim
    \upperfontdim2\fontdim
    \fontdim1.5\fontdim
  \else
    \loop
      \fontdim2\fontdim % Double font size
      \buildbox
    \ifdim\dimen@<#2 \repeat
    \upperfontdim\fontdim
    \lowerfontdim.5\fontdim
    \fontdim.75\fontdim
  \fi
  % Now try to find the optimum size
  \loop
    %\message{Bounds: \the\lowerfontdim\space
    %         \the\fontdim\space \the\upperfontdim^^J}
    \buildbox
    \ifdim\dimen@>#2
      \moreiterationstrue
      \upperfontdim\fontdim
      \advance\fontdim\lowerfontdim
      \fontdim.5\fontdim
    \else
      \advance\dimen@-#2
      \ifdim\dimen@<10pt
        \lowerfontdim\fontdim
        \advance\fontdim\upperfontdim
        \fontdim.5\fontdim
        \dimen@\upperfontdim
        \advance\dimen@-\lowerfontdim
        \ifdim\dimen@<.2pt
          \moreiterationsfalse
        \else
          \moreiterationstrue
        \fi
      \else
        \moreiterationsfalse
      \fi
    \fi
  \ifmoreiterations \repeat
  \box0% Typeset content
}
\makeatother

% ---------------------------------------------------------------------------
% IDENTIFICAÇÃO
% ---------------------------------------------------------------------------
\titulo{Análise do impacto de variáveis de Recuperação de Informação no desempenho de classificadores de Mineração de Textos}
\autor{Ruan de Medeiros Bahia}
\local{Juazeiro - BA}
\orientador{Prof. Dr. Rosalvo Ferreira de Oliveira Neto}
%\coorientador{M. Sc. Ciclano Fulado de Tal}
\instituicao{
FUNDAÇÃO UNIVERSIDADE FEDERAL DO VALE DO SÃO FRANCISCO
	\\
COLEGIADO DE ENGENHARIA DE COMPUTAÇÃO}
\tipotrabalho{Formulário de matrícula de TCC I}
\preambulo{Trabalho apresentado à Universidade Federal do Vale do São Francisco - Univasf, Campus Juazeiro, como requisito da obtenção do título de Bacharel em Engenharia de Computação.}


% -----------------------------------------------------------------------------
% CONFIGURACAO DO SUMARIO - by @Gabrielr2508
% Precisa estar aqui, por isso não foi para o commands.tex, não descobrimos o motivo, %caso saiba, por favor, faça um pull request! :D
% -----------------------------------------------------------------------------
% Secao primaria (Chapter) Caixa alta, Negrito, tamanho 12
\makeatletter
\renewcommand*{\l@chapter}[2]{%
  \l@chapapp{\uppercase{#1}}{#2}{\cftchaptername}}
\makeatother
% Secao secundaria (Section) Caixa baixa, Negrito, tamanho 12
\renewcommand{\cftsectionfont}{\uppercase} %ponha \rmfamily se quiser serifadas...

% Secao terciaria (Subsection) Caixa baixa, negrito, tamanho 12
\renewcommand{\cftsubsectionfont}{\bfseries}

% Secao quaternaria (Subsubsection) Caixa baixa, tamanho 12
\renewcommand{\cftsubsubsectionfont}{\normalfont}

% Seção quinaria (subsubsubsection) Caixa baixa, sem negrito, tamanho 12
\renewcommand{\cftparagraphfont}{\normalfont\itshape}

% -----------------------------------------------------------------------------
% Início do TCC 
% -----------------------------------------------------------------------------
\begin{document}
\selectlanguage{brazil}
	\frenchspacing % Retira espaço extra obsoleto entre as frases.
	
    \setlength{\belowcaptionskip}{0pt}
    \setlength{\abovecaptionskip}{0pt}
    \setlength{\intextsep}{-18pt}
        \begin{figure}[h]
    		\begin{center}
    		   \includegraphics[scale=0.8]{img/brasao_brasil_m.png}
    		\end{center}
    	\end{figure}
        \center
        %\textscale{0.9} 
    	{\textbf{\ABNTEXchapterfont\imprimirinstituicao}}
    	\\
        Avenida Antonio Carlos Magalhães, 510 - Santo Antônio CEP: 48902-300 \\
        Juazeiro/BA, Tel/Fax: (74) 2102-7636 - Juazeiro - BA \\
        \url{www.univasf.edu.br/~ccomp}
        
    	\begin{mdframed}[roundcorner=10pt]
    	    \begin{minipage}[t][3\height][c]{\dimexpr\textwidth-1\fboxsep-2\fboxrule\relax}
                \centering
                \ABNTEXchapterfont\large\imprimirtipotrabalho
            \end{minipage}
        \end{mdframed}
        
        
        \newmdenv[%
            roundcorner=5pt,
            subtitlebelowline=true,subtitleaboveline=true,
            subtitlebackgroundcolor=yellow!70!white, %             backgroundcolor=blue!20!white,
            frametitle={Título do projeto de pesquisa},
            frametitlerule=true, 
            frametitlebackgroundcolor=yellow!70!white,
            subtitlebelowskip=0.1\baselineskip,
            subtitleaboveskip=0.5\baselineskip,
            subtitleinneraboveskip=5pt,
            subtitleinnerbelowskip=5pt,
            ]{subtitleenv}
        \begin{subtitleenv}
            \imprimirtitulo

            \mdfsubtitle{Nome do discente}
            {\small
            \imprimirautor
            }
            \mdfsubtitle{Nome do orientador e co-orientador}
            {\small
            \imprimirorientador
            }
            \mdfsubtitle{Resumo e palavras-chave}
            {\small
            %--------------------------------------------------------------------------------------
% Resumo e palavras-chave
%--------------------------------------------------------------------------------------
O mundo informatizado gera uma quantidade gigantesca de dados textuais diariamente, e a Mineração de Texto objetiva transformar esses dados em informações úteis, em conhecimento, tendo aplicação inclusive na área forense.
A área de Recuperação de Informação contribui para o desenvolvimento da Mineração de Texto no pré-processamento, porém sua utilização direta na criação de atributos não é usual, sendo proposto pela primeira vez por \citeonline{WEREN_CLEF_2014}.

A partir de uma revisão bibliográfica das áreas de Recuperação de Informação e de Mineração de Texto, o autor sugere uma metodologia para criação de atributos utilizando a função de ranqueamento BM25, similar à utilizada por \citeonline{WEREN_MESTRADO_2014}, e utilizando ferramentas de armazenamento e indexação já existentes para o cálculo.
Nesta metodologia, a análise do desempenho dos novos atributos sugeridos é feita por meio da mensuração do desempenho de classificador por medidas consolidadas na literatura de Mineração de Texto.


Palavras-chave: Mineração de Texto, Recuperação de Informação, criação de atributos, avaliação de desempenho, engenharia de atributos.
            }
            \mdfsubtitle{Qualificação do problema a ser abordado}
            {\small
            %--------------------------------------------------------------------------------------
% Qualificação do problema a ser abordado 
%--------------------------------------------------------------------------------------

A Mineração de Textos consiste na extração de informações de interesse e de padrões não triviais em documentos de texto não-estruturados \cite{Srivastava2009TMCCA}. 
Para tanto, nesse processo de descoberta de conhecimento, são utilizadas diversas técnicas de inteligência artificial, principalmente de seus ramos pais, a mineração de dados e o aprendizado de máquina, visionando um bom proveito dessa descoberta de conhecimento.

Além dos campos citados acima, o campo da Mineração de Textos se utiliza de métodos de processamento de linguagem natural, aprendizagem estatística, web semântica, e Recuperação de Informação \cite{Sammut2017EMLDM} \cite{Manning2008IIR}. 
Dentre os passos da Mineração de Textos para extração de conhecimento de um conjunto de dados, a seleção do conjunto de parâmetros a ser utilizado por um classificador se faz bastante importante pois estes parâmetros influem diretamente no desempenho dos mesmos.

            }
            \mdfsubtitle{Justificativa}
            {\small
            %--------------------------------------------------------------------------------------
% Justificativa
%--------------------------------------------------------------------------------------

A Mineração de Textos proveem diversas técnicas a fim de extrair conhecimento de grandes volumes de dados, oferecendo principalmente técnicas de clusterização de dados, visualização e classificação \cite{Sammut2017EMLDM}. 
Dentre as aplicação práticas da mineração de texto temos que um dos grandes ramos de interesse é análise de autoria \cite{Berry2010TMAT}, que pode servir tanto para identificar casos de plágio como combater crimes cibernéticos, identificando predadores sexuais que tentam se passar por outras pessoas, tendo assim grande relevância como ferramenta de auxílio a esses casos.

Melhorar o desempenho dos classificadores de Mineração de Texto é uma tarefa desafiadora e de grande interesse prático para a ajuda no combate a esses crimes\cite{Berry2010TMAT}, além também da melhora nestes classificadores ajudar diretamente em qualquer tarefa de classificação de Mineração de Textos.
            }
            \mdfsubtitle{Objetivo a ser alcançado}
            {\small
            %--------------------------------------------------------------------------------------
% Objetivo a ser alcançado 
%--------------------------------------------------------------------------------------

O objetivo geral deste trabalho é fazer uma avaliação do desempenho de atributos oriundos de Recuperação de Informação para tarefas de Mineração de Textos.
Ainda, possui os seguintes objetivos específicos:
\begin{itemize}
    \item Avaliar o ganho de desempenho de classificadores de Mineração de Texto com adição de atributos derivados da função de ranqueamento BM25 da Recuperação de Informação, em pelo menos 2 corpus de competições diferentes, utilizando medidas consolidadas na literatura;
    
    \item Reproduzir soluções disponíveis \textit{online} para os corpus selecionados, comprovando as medidas dos resultados das competições;
    
    \item Elencar em qual dos corpus selecionados os atributos criados proporcionam maior ganho de desempenho de classificador;
    
    \item Comparar o desempenho computacional de ferramentas de armazenamento e indexação de textos:
    \begin{itemize}
        \item na questão de indexação;
        \item na questão de consulta utilizando as implementações do BM25 nativas das ferramentas;
    \end{itemize}
    
    \item Avaliar, empiricamente, a facilidade de instalação, utilização e integração das ferramentas de armazenamento e indexação selecionadas.
\end{itemize}
            }
            \mdfsubtitle{Metodologia a ser empregada}
            {\small
            %--------------------------------------------------------------------------------------
% Metodologia a ser empregada 
%--------------------------------------------------------------------------------------
Levantamento bibliográfico, feito com auxílio do Google Scholar \url{https://scholar.google.com.br/} para a pesquisa inicial, sobre Recuperação de Informação, Mineração de Textos, Mineração de Dados, Análise de Textos e Classificação de Autores.
Será utilizado como critério de seleção dos artigos e livros:
\begin{itemize}
    \item Estar escrito em português ou inglês; e
    \item Ter no resumo/sumário quaisquer das palavras-chave: Recuperação de Informação; Mineração de Textos; Mineração de Dados, Análise de Textos; Classificação de Autores; Information Retrieval; Text Mining; Data Mining; Text Analysis; e Autor Profiling; e
    \item Disponibilidade gratuita pela internet, podendo ser encontrado utilizando o Google ou Google Scholar ou disponibilizado pela Library Genesis \url{http://gen.lib.rus.ec/} ou Sci-Hub \url{http://sci-hub.tw/}.
\end{itemize}
Haverá priorização pelos artigos e livros mais recentes, e também pelos que forem mais citados em outros artigos que atendam ao tema.

Os artigos e livros serão categorizados de acordo com as cinco palavras chave estabelecidas acima nos critérios do levantamento, sendo então feita uma revisão livre dos artigos escolhidos como relevantes pelo autor, sendo estes lidos em sua totalidade, anotados e resumidos e devidamente fichados. E os livros encontrados serão inicialmente analisados pelo seu sumário, sendo localizados os capítulos pertinentes aos temas do levantamento, e estes então lidos e também resumidos.

O processo de levantamento bibliográfico será feito interativamente com a escrita do trabalho, e este último também será desenvolvido em paralelo com o processo de estudo posterior das técnicas e ferramentas de RI (Recuperação de Informação) a serem utilizadas como parâmetro em classificadores de MT (Mineração de Texto).

A seleção dos banco de dados a serem aplicados os classificadores de MT será feita após ter sido concluída a revisão bibliográfica e terem sido já avaliados quais grupos de parâmetros, ou variáveis, de RI serão utilizados e quais os respectivos classificadores de MT que vão utilizar dessas variáveis.
            }
            \mdfsubtitle{Referencial bibliográfico}
            {\small
            %--------------------------------------------------------------------------------------
% Referências bibliográficas
%--------------------------------------------------------------------------------------
% Inserindo a bibliografia sem o título
\begingroup
    \renewcommand{\section}[2]{}%
    \renewcommand{\chapter}[2]{}% for other classes
    \bibliography{tex/references}
\endgroup

            }
            \mdfsubtitle{Cronograma de atividades}
            {\small
            %--------------------------------------------------------------------------------------
% Cronograma de atividades
%--------------------------------------------------------------------------------------

%\section{Cronograma} \label{sec:Cronograma}

%A tabela \ref{tab:cronograma} mostra o cronograma de atividades a serem executadas para o TCC II, com base no calendário de 201X.Y da UNIVASF.

%\newpage
%\begin{table}[!thb]
    %\huge
\begin{center}
    \captionof{table}{\label{tab:cronograma} Cronograma das atividades previstas para o TCC II}
    \begin{adjustbox}{max width={\textwidth},keepaspectratio}%
    \begin{tabular}{@{} p{6.5cm}|c|c|c|c|c|c @{}}
        \toprule
        \textbf{Atividade}
        & Set & Out & Nov & Dez & Jan & Fev
        \\ \hline
        Definição e obtenção dos corpus para avaliação 
        & X   &     &     &     &     &          
        \\ \hline
        Inspeção e seleção das soluções com código fonte disponível
        & X   &     &     &     &     &          
        \\ \hline
        Instalação e familiarização com as ferramentas de arquivamento e indexação
        & X   & X   &     &     &     &          
        \\ \hline
        Indexação do conjunto de treinamento dos corpus
        &     & X   & X   &     &     &          
        \\ \hline
        Adição dos atributos de RI às soluções selecionadas
        &     & X   & X   &     &     &    
        \\ \hline
        Mineração dos dados por meio da reprodução das soluções selecionadas com/sem adição dos atributos de RI
        &     &     & X   & X   & X   &  
        \\ \hline
        Escrita do TCC II                       
        & X   & X   & X   & X   & X   &         
        \\ \hline
        Defesa do TCC II                        
        &     &     &     &     & X   &        
        \\
        \bottomrule
    \end{tabular}
    \end{adjustbox}
    \legend{\ABNTEXfontereduzida \textbf{Fonte:} O autor.}
\end{center}
%\end{table}
            }
            \mdfsubtitle{Necessidades e disponibilidade de recursos e infraestrutura para o desenvolvimento deste projeto}
            {\small
            %--------------------------------------------------------------------------------------
% Necessidades e disponibilidade de recursos e infraestrutura para o desenvolvimento deste projeto  
%--------------------------------------------------------------------------------------
            }
           
        \end{subtitleenv}
        
        
        \vspace*{1cm}
        \hrulefill \\
        \imprimirorientador \\
        Orientador(a) \\
        
        % \vspace*{1cm}
        % \hrulefill \\
        % NOME DO CO-ORIENTADOR \\
        % Co-Orientador(a) \\
        
        \vspace*{1cm}
        \hrulefill \\
        \imprimirautor \\
        Aluno(a) \\
        
        \vspace*{1cm}
        \imprimirlocal, 13/05/2019.\\

\end{document}
